\documentclass[12pt,letterpaper]{article}
\usepackage[spanish,es-noshorthands]{babel}
\usepackage[utf8]{inputenc}
\usepackage[T1]{fontenc}
\usepackage[margin=1in]{geometry}
\usepackage{newtxtext,newtxmath}
\usepackage{setspace}
\usepackage{fancyhdr}
\usepackage{titlesec}
\usepackage{amsmath,amssymb}
\usepackage{graphicx}
\usepackage[hidelinks]{hyperref}
\usepackage{url}
\doublespacing
\setlength{\parindent}{0.5in}
\setlength{\parskip}{0pt}
\setcounter{secnumdepth}{0}
\titleformat{\section}{\normalfont\bfseries\centering}{}{0pt}{}
\titleformat{\subsection}{\normalfont\bfseries}{}{0pt}{}
\pagestyle{fancy}
\fancyhf{}
\fancyhead[R]{\thepage}
\renewcommand{\headrulewidth}{0pt}
\begin{document}
\begin{titlepage}
\thispagestyle{fancy}
\centering
\vspace*{0.8in}
{\bfseries Modelo de optimización para la reprogramación de parciales universitarios\par}
\vspace{1\baselineskip}
{Manuel David Polo Salgado\par
Andrés David Issa Camargo\par
Daniel Alejandro Rincón Navarro\par
Cristian Camilo Fuentes Sánchez\par
Matías Guarín Rueda\par}
\vspace{1\baselineskip}
{Universidad del Norte\par
MAT 4021: Matemáticas Discretas\par
Profesor: MSc. Iván Gómez\par
16 de noviembre de 2026\par}
\end{titlepage}
\section*{Resumen}
Este trabajo describe CronoGraph, una aplicación implementada para apoyar la programación de parciales universitarios. El problema se formula mediante un grafo en el que cada nodo representa una evaluación y las aristas representan conflictos entre parciales que comparten estudiantes. El modelo pondera esos conflictos según la cantidad de estudiantes afectados, la proximidad temporal de las evaluaciones y la importancia académica de cada parcial. Esta importancia se estima a partir de los créditos de la asignatura, el porcentaje de la nota final correspondiente al parcial y el índice histórico de retiro. Se describen los criterios, las alternativas consideradas y la función objetivo para comparar calendarios factibles. El simulador implementa búsqueda local con reinicios aleatorios y el solucionador CP-SAT de Google OR-Tools; también permite ejecutarlos en modo combinado. Una prueba de tres parciales contrasta el resultado de la búsqueda con la enumeración de las ocho asignaciones posibles; quedan pendientes experimentos con datos institucionales.


\noindent\textit{Palabras clave:} 
programación de exámenes, optimización, teoría de grafos, calendario académico, parciales


\section{Introducción: problema, impacto y objetivos}
La programación de parciales puede concentrar evaluaciones en periodos cortos y afectar simultáneamente a muchos estudiantes, un problema estudiado en la literatura de horarios de exámenes (Carter et al., 1994; Qu et al., 2009; Bassimir y Wanka, 2025). Además de evitar coincidencias de horario, se busca reducir la proximidad entre evaluaciones compartidas, respetando las fechas y horarios en que cada parcial puede realizarse. El proyecto aborda este problema mediante una aplicación que cuantifica los conflictos, valida calendarios y compara propuestas de reprogramación. El código fuente y las instrucciones de ejecución están disponibles en \url{https://github.com/WosDan/CronoGraph}.

Se escogió este problema por su relación directa con la organización de la carga académica y por la posibilidad de representar sus componentes con datos observables: estudiantes inscritos en cada materia, fechas disponibles, créditos, ponderación de los parciales e historial de retiros. Un calendario que considere simultáneamente estos factores puede ofrecer recomendaciones más informadas que una asignación basada únicamente en evitar choques de horario.

\section{Modelación con la teoría del curso}
El problema se representa mediante un grafo ponderado. Cada nodo corresponde a un parcial y cada arista conecta dos parciales cuando hay estudiantes que deben presentar ambos. El simulador considera para cada evaluación su conjunto de fechas y horarios permitidos, construye o utiliza la relación de estudiantes compartidos y evalúa las combinaciones de fechas con una función de costo. La disponibilidad se trata como una restricción obligatoria; la reducción de la proximidad y de la concentración de evaluaciones se trata como criterio de calidad.

\subsection{Disponibilidad y criterios del modelo}
Para cada parcial $i$, se define un conjunto $D_i$ de fechas y horarios autorizados. Este conjunto puede derivarse del horario habitual de clase cuando el parcial se realiza durante la sesión, o de las opciones especiales aprobadas por el docente o la institución. El calendario asignado debe seleccionar una opción perteneciente a $D_i$.

Además, la fecha propuesta puede adelantarse o aplazarse como máximo 14 días calendario respecto de la fecha inicial, incluidos los extremos. Las franjas autorizadas deben reflejar la disponibilidad real del docente y el avance de los contenidos: el simulador no verifica automáticamente que los temas evaluados ya se hayan enseñado. Los solapamientos de estudiantes, docentes y aulas se comprueban por separado.

Los criterios del modelo son:
\begin{itemize}
  \item \textbf{Disponibilidad:} cada parcial solo puede asignarse a las fechas y horas autorizadas.
  \item \textbf{Estudiantes compartidos:} dos evaluaciones que afectan a más estudiantes generan un conflicto potencial mayor.
  \item \textbf{Proximidad temporal:} se penalizan los parciales separados por cero, uno, dos o tres días calendario; a partir de cuatro días la penalización es nula.
  \item \textbf{Importancia académica:} el peso del parcial combina créditos, porcentaje de la nota final e índice histórico de retiro.
\end{itemize}

\subsection{Peso académico y función objetivo}
Los créditos representan la carga académica formal prevista para una asignatura (Presidencia de la República de Colombia, 2019) y el porcentaje del parcial representa su contribución a la calificación final. Se combinan multiplicativamente para que el indicador dependa de ambas dimensiones. El índice histórico de retiro se incorpora como factor complementario de priorización, no como prueba causal de dificultad o de que una fecha provoque retiros. La fórmula y el ajuste del 50\% son decisiones de diseño del proyecto, susceptibles de validación.

El peso académico del parcial $i$ se define como
\begin{equation}
  W_i = C_i p_i\left(1+\frac{r_i}{2}\right),
\end{equation}
donde $C_i$ es el número de créditos de la asignatura, $p_i$ es la fracción de la nota final correspondiente al parcial y $r_i\in[0,1]$ es el índice histórico de retiro. Así, el factor de retiro aumenta el peso base entre 0\% y 50\%.

Para dos parciales $i$ y $j$, la penalización por proximidad, con $\Delta_{ij}$ como la distancia en días calendario entre sus fechas, es
\begin{equation}
  P(\Delta_{ij}) = \max(0,4-\Delta_{ij}).
\end{equation}
Si $A_{ij}$ es el número de estudiantes que presentan ambos parciales, el conflicto ponderado se calcula como
\begin{equation}
  K_{ij} = A_{ij}P(\Delta_{ij})\frac{W_i+W_j}{2}.
\end{equation}
La función objetivo del calendario es minimizar la suma de los conflictos entre pares:
\begin{equation}
  F = \sum_{i<j} K_{ij}.
\end{equation}
Una extensión posible, que no forma parte de la función objetivo implementada, es
\begin{equation}
  \min F^* = \sum_{i<j}K_{ij} + \lambda\sum_i|d_i-d_i^0| + \mu B_{\max},
\end{equation}
donde $d_i$ es la fecha nueva, $d_i^0$ la fecha original y $B_{\max}$ la mayor concentración de parciales que afronta un estudiante en una ventana definida. $\lambda$ y $\mu$ representan la importancia relativa de conservar el calendario y reducir esa concentración.

\subsection{Centralidades en el grafo de conflictos}
El grafo de conflictos permite calcular el grado ponderado, que suma la intensidad de las conexiones de cada parcial, y la intermediación, que mide su presencia en caminos mínimos. Para esta última se utiliza la distancia $1/H_{ij}$, donde $H_{ij}$ es la intensidad del conflicto. Ambas medidas orientan el orden de construcción de algunas soluciones iniciales; en la búsqueda local, el grado ponderado y después la intermediación se usan para desempatar movimientos con mejoras prácticamente iguales. La cercanía se presenta como indicador, pero no selecciona movimientos. El efecto de esta priorización debe medirse comparando ejecuciones con y sin centralidades.

\section{Solución: algoritmos y datos}
El conjunto de demostración incluido en el repositorio contiene 120 estudiantes ficticios, 15 asignaturas, 12 docentes, 20 grupos y 20 parciales. Sus datos fueron generados con semilla 42; hay 595 matrículas y cada estudiante cursa entre cuatro y seis asignaturas. El usuario también puede generar otra instancia sintética o cargar un archivo JSON personalizado que supere la validación. Los datos incluidos son sintéticos y no representan registros institucionales reales.

El backend verifica las relaciones entre entidades, las franjas autorizadas, la duración del examen y los solapamientos de estudiantes, docentes y aulas. Las franjas permitidas son sesiones habituales de clase o espacios sabatinos marcados explícitamente como autorizados. A partir de la matrícula forma la matriz binaria estudiante--parcial $B$ y la matriz $A=B^{\mathsf T}B$, cuyos elementos fuera de la diagonal cuentan los estudiantes compartidos. Con estas cantidades construye el grafo ponderado y evalúa el calendario inicial y el propuesto.

El límite de reprogramación de cada evaluación es $|d_i-d_i^0|\leq 14$ días calendario, además de pertenecer a su dominio autorizado. El calendario permite examinar la franja asignada y la factibilidad global. La sección de estudiantes compara sus pares de parciales separados por hasta tres días; una reducción de este indicador no demuestra por sí sola una mejora del bienestar o del rendimiento académico.

La entrada incluye un conjunto JSON con estudiantes, matrículas, asignaturas, docentes, parciales y franjas autorizadas. Como salida, la interfaz presenta el calendario propuesto, su factibilidad, costos inicial y final, reducción, cambios de horario, métricas por estudiante, distribución de penalizaciones y traza de convergencia. La interfaz permite explorar datos y grafos y elegir búsqueda local, CP-SAT o el modo combinado. Los resultados de optimización pueden guardarse como archivos JSON en el backend y exportarse desde la aplicación.

\subsection{Métodos de optimización}
Además de las alternativas conceptuales para estimar carga y ampliar la función objetivo, la implementación compara dos métodos de optimización. La búsqueda local con reinicios aleatorios parte del calendario inicial si es factible y construye otros calendarios mediante una heurística voraz aleatorizada. En cada reinicio realiza búsqueda de mejor mejora hasta un óptimo local. Si el calendario inicial es factible, el sistema garantiza que el costo retornado no sea mayor que el inicial.

La segunda alternativa es CP-SAT de Google OR-Tools (Google, s. f.). El modelo selecciona exactamente una franja válida por parcial, prohíbe solapamientos incompatibles y minimiza el costo temporal con coeficientes enteros escalados por $10^4$. El límite de tiempo es configurable; solo el estado \texttt{OPTIMAL} certifica el óptimo del objetivo escalado. El modo combinado proporciona a CP-SAT la solución de búsqueda local como punto de partida cuando es factible y presenta la solución factible de menor costo original entre ambas. Esta es la estrategia comparativa que efectivamente ofrece la aplicación.

\subsection{Alternativas y extensiones}
Se consideran como alternativas las horas reales de preparación registradas por estudiantes, las horas estimadas por docentes, y los resultados históricos de los parciales (por ejemplo, promedios o tasas de reprobación). Las horas reales serían una medida más directa del esfuerzo temporal, pero requerirían encuestas de seguimiento o diarios de estudio. Las calificaciones históricas serían complementarias, pues un resultado bajo no demuestra por sí mismo que se haya dedicado más tiempo. Se propone comparar los indicadores con las horas de estudio registradas y evaluar su asociación mediante correlación de rangos de Spearman (National Institute of Standards and Technology [NIST], s. f.; United Nations Statistics Division, s. f.).

Una extensión podría penalizar también los cambios respecto del calendario original y la máxima cantidad de parciales que afronta un estudiante en una ventana temporal. Esto permitiría equilibrar la reducción de conflictos con la estabilidad del calendario y la carga individual; los coeficientes de ponderación tendrían que calibrarse.

Como alternativas futuras se pueden estudiar recocido simulado y algoritmos genéticos, que exploran otras asignaciones factibles, así como formulaciones de programación entera. Se propone compararlos con los métodos implementados usando las mismas instancias, función objetivo y restricciones, registrando costo, factibilidad, tiempo de ejecución y variabilidad entre ejecuciones. Estas alternativas se proponen para evaluación futura; no se afirma que ya hayan sido probadas.

El calendario de menor costo entre los algoritmos evaluados es el mejor encontrado en esos experimentos, pero no necesariamente el mejor de todos los calendarios posibles. Para instancias pequeñas puede verificarse el mínimo enumerando todas las asignaciones factibles. En instancias mayores puede recurrirse a un método exacto hasta obtener un certificado de optimalidad; para que esa garantía corresponda a la función original debe evitarse la pérdida de precisión de los coeficientes, siempre que la representación entera sea viable. Los recursos y el tiempo disponibles pueden impedir completar la demostración.

\section{Resultados e interpretación}
Para validar la implementación se utilizó una instancia sintética pequeña de tres parciales y cuatro estudiantes, con dos franjas por parcial. La enumeración de las ocho combinaciones posibles encontró que las ocho eran factibles y que el costo mínimo era $F=0$. El calendario inicial tenía costo $F=23.14875$. En la ejecución de búsqueda local con diez reinicios y semilla 42 se obtuvo $F=0$, moviendo dos de los tres parciales; por tanto, para esta instancia concreta el resultado coincide con el mínimo comprobado por enumeración. Este caso no demuestra optimalidad en otras instancias ni constituye una medición de bienestar estudiantil.

En esta revisión también se ejecutaron ocho pruebas seleccionadas de restricciones y optimización: todas pasaron. Incluyen la comparación contra enumeración en la instancia pequeña y el límite de reprogramación de 14 días. Para evaluar el rendimiento con la instancia de 20 parciales o con datos institucionales, aún deben registrarse ejecuciones, tiempos, costos y resultados comparables.

\section{Conclusiones y limitaciones}
CronoGraph implementa la reprogramación de parciales mediante un grafo ponderado, restricciones de factibilidad y una función objetivo que combina estudiantes afectados, proximidad temporal e importancia académica. La aplicación permite comparar búsqueda local con reinicios aleatorios y CP-SAT. La calidad de las recomendaciones todavía debe evaluarse cuantitativamente en ejecuciones documentadas y mediante validación de los parámetros con datos institucionales.

El modelo depende de la calidad y disponibilidad de los datos de matrícula, retiro, créditos, ponderación y fechas autorizadas. El índice de retiro puede depender de factores personales, económicos, docentes y académicos, por lo que no debe interpretarse causalmente. El límite de cuatro días y el incremento máximo del 50\% son parámetros propuestos por el proyecto y necesitan validación empírica. La búsqueda local puede terminar en un óptimo local, y la heurística voraz no garantiza encontrar una asignación factible si existe. CP-SAT puede alcanzar el límite de tiempo sin probar optimalidad; además, la certificación usa coeficientes redondeados a escala $10^4$. Los resultados del caso pequeño no permiten extrapolar el desempeño del sistema a instancias institucionales reales.

Como trabajo futuro se propone validar el indicador de peso frente a horas efectivas de estudio y calibrar sus parámetros mediante análisis de sensibilidad con actores académicos. El sistema ya implementa restricciones de solapamiento de estudiantes, docentes y aulas, y presenta métricas de concentración e impacto individual. Falta evaluar experimentalmente esas métricas con instancias representativas y cuantificar la reducción de conflictos, los estudiantes beneficiados y los parciales reprogramados.

\section*{Referencias}
\begingroup
\setlength{\parindent}{0pt}
\setlength{\hangindent}{0.5in}
Bassimir, B., y Wanka, R. (2025). On the computation of robust examination timetables: Methods and experimental results. \textit{Journal of Scheduling, 28}, 159--181. \url{https://doi.org/10.1007/s10951-024-00815-y}

\hangindent=0.5in
Carter, M. W., Laporte, G., y Chinneck, J. W. (1994). A general examination scheduling system. \textit{Interfaces, 24}(3), 109--120. \url{https://doi.org/10.1287/inte.24.3.109}

\hangindent=0.5in
Google. (s. f.). \textit{CP-SAT solver}. OR-Tools. \url{https://developers.google.com/optimization/cp/cp_solver}

\hangindent=0.5in
National Institute of Standards and Technology. (s. f.). \textit{Rank correlation}. \url{https://www.itl.nist.gov/div898/software/dataplot/refman2/ch2/rankcorr.pdf}

\hangindent=0.5in
Presidencia de la República de Colombia. (2019). \textit{Decreto 1330 de 2019}. \url{https://www.funcionpublica.gov.co/eva/gestornormativo/norma.php?i=98270}

\hangindent=0.5in
Qu, R., Burke, E. K., McCollum, B., Merlot, L. T. G., y Lee, S. Y. (2009). A survey of search methodologies and automated system development for examination timetabling. \textit{Journal of Scheduling, 12}, 55--89.

\hangindent=0.5in
United Nations Statistics Division. (s. f.). \textit{Time-use diaries}. \url{https://unstats.un.org/UNSDWebsite/demographic-social/time-use/guidelines/III/}
\endgroup
\end{document}
